\section{Methodology}\label{sec:study_design}
\begin{figure*}[t]
\centering
\resizebox{\textwidth}{!}{%
\begin{tikzpicture}[
  x=1cm,y=1cm,
  font=\sffamily,
  >={Stealth[length=2.3mm,width=1.7mm]},
  panel/.style n args={2}{rounded corners=3.5mm, line width=0.8pt, draw=#1, fill=#2, outer sep=0pt},
  card/.style={rounded corners=2.6mm, line width=0.6pt, draw=#1, fill=white, outer sep=0pt},
  head/.style={font=\sffamily\bfseries\fontsize{12.5}{14}\selectfont,text=ink,anchor=west},
  sub/.style={font=\sffamily\fontsize{9}{11}\selectfont,text=subink!55,anchor=west},
  section/.style={font=\sffamily\bfseries\fontsize{9}{10}\selectfont,anchor=west},
  bigtitle/.style={font=\sffamily\bfseries\fontsize{15.5}{17}\selectfont,text=ink,anchor=west},
  iconbox/.style={rounded corners=2.2mm,minimum width=0.92cm,minimum height=0.92cm,inner sep=0pt,outer sep=0pt},
  flow/.style={-{Stealth[length=2.1mm]},line width=1.05pt},
  thickflow/.style={-{Stealth[length=2.4mm]},line width=1.35pt}
]
% Canvas hugs the panels so the scaled figure is not wider than the artwork
\fill[paper] (-0.22,0.18) rectangle (27.42,17.82);

% ===== 01 Discovery panel =====
\node[panel={tealLine}{tealBg},minimum width=16.9cm,minimum height=6.45cm,anchor=north west] (discover) at (0,17.6) {};
% header badge and title
\fill[teal] (1.02,16.45) circle (0.62);
\node[white,font=\bfseries\fontsize{16}{16}\selectfont] at (1.02,16.45) {01};
\node[section,text=teal] at (2.05,16.98) {\textls[150]{DISCOVER}};
\node[bigtitle] at (2.05,16.48) {Repository discovery \& screening};

% Source discovery card
\node[card={tealLine},minimum width=8.35cm,minimum height=1.9cm,anchor=north west] (source) at (0.4,15.25) {};
\node[iconbox,fill=teal!14] at (1.35,14.3) {};
\node[text=teal!70!black,font=\fontsize{21}{21}\selectfont] at (1.35,14.3) {\faSearch};
\node[head] at (2.45,14.62) {Source discovery};
\node[sub] at (2.45,14.12) {GitHub search};

% Eligibility card
\node[card={tealLine},minimum width=6.45cm,minimum height=1.9cm,anchor=north west] (elig) at (9.75,15.25) {};
\node[iconbox,fill=teal!14] at (10.7,14.3) {};
\node[text=teal!75!black,font=\fontsize{21}{21}\selectfont] at (10.7,14.3) {\faFilter};
\node[head] at (11.8,14.62) {Eligibility screening};
\node[sub] at (11.8,14.12) {4,743 passed};

\draw[thickflow,teal!80!black] (8.88,14.3) -- (9.62,14.3);

% Out of scope
\draw[flow,red] (13.0,13.35) -- (13.0,12.72);
\node[card={red},fill=redBg,minimum width=5.75cm,minimum height=1.35cm,anchor=north west] (scope) at (10.45,12.7) {};
\node[iconbox,fill=red!12,minimum width=0.82cm,minimum height=0.82cm] at (11.2,12.02) {};
\node[text=red,font=\fontsize{16}{16}\selectfont] at (11.2,12.02) {\faMinusCircle};
\node[font=\bfseries\fontsize{9.5}{10}\selectfont,text=red!85!black,anchor=west] at (12.0,12.2) {Out of scope};
\node[font=\fontsize{8.4}{10}\selectfont,text=subink!55,anchor=west] at (12.0,11.83) {Templates or $<6$ months};
\node[font=\bfseries\fontsize{9}{10}\selectfont,text=subink!75,anchor=east] at (15.9,12.2) {$\sim$10,799};

% ===== 03 Structure panel =====
\node[panel={purpleLine}{purpleBg},minimum width=9.15cm,minimum height=8.55cm,anchor=north west] (structure) at (17.55,17.6) {};
\fill[purple] (18.45,16.45) circle (0.62);
\node[white,font=\bfseries\fontsize{16}{16}\selectfont] at (18.45,16.45) {03};
\node[section,text=purple] at (19.45,16.98) {\textls[150]{STRUCTURE}};
\node[bigtitle] at (19.15,16.48) {Extract rules \& history};

% Structure cards
\node[card={purpleLine},minimum width=7.45cm,minimum height=1.52cm,anchor=north west] (extract) at (17.9,15.55) {};
\node[iconbox,fill=purple!11,minimum width=1.0cm,minimum height=1.0cm] at (18.85,14.79) {};
\node[text=purple,font=\fontsize{20}{20}\selectfont] at (18.85,14.79) {\faFile*[regular]};
\node[head] at (19.85,15.02) {Rule extraction};
\node[sub] at (19.85,14.54) {Parse 11,136 rule files};

\node[card={purpleLine},minimum width=7.45cm,minimum height=1.52cm,anchor=north west] (normalize) at (17.9,13.58) {};
\node[iconbox,fill=purple!11,minimum width=1.0cm,minimum height=1.0cm] at (18.85,12.82) {};
\node[text=purple,font=\fontsize{20}{20}\selectfont] at (18.85,12.82) {\faFilter};
\node[head] at (19.85,13.05) {Rule normalization};
\node[sub] at (19.85,12.57) {Constructs $\rightarrow$ One schema};

\node[card={purpleLine},minimum width=7.45cm,minimum height=1.52cm,anchor=north west] (taxonomy) at (17.9,11.61) {};
\node[iconbox,fill=purple!11,minimum width=1.0cm,minimum height=1.0cm] at (18.85,10.85) {};
\node[text=purple,font=\fontsize{20}{20}\selectfont] at (18.85,10.85) {\faSitemap};
\node[head] at (19.85,11.08) {Taxonomy construction};
\node[sub] at (19.85,10.60) {Rule categories};

\draw[flow,blue!75!black] (extract.south) -- (normalize.north);
\draw[flow,blue!75!black] (normalize.south) -- (taxonomy.north);

% ===== Lower panels =====
% 02 workflow integration
\node[panel={orangeLine}{orangeBg},minimum width=8.35cm,minimum height=8.25cm,anchor=north west] (workflow) at (0,8.65) {};
\fill[orange] (1.05,7.65) circle (0.62);
\node[white,font=\bfseries\fontsize{16}{16}\selectfont] at (1.05,7.65) {02};
\node[section,text=orange!85!black] at (2.05,8.05) {\textls[150]{DETECT}};
\node[bigtitle] at (2.05,7.58) {Workflow integration};

\node[card={orangeLine},minimum width=7.55cm,minimum height=1.65cm,anchor=north west] (wfscan) at (0.4,6.65) {};
\node[iconbox,fill=orange!15] at (1.35,5.83) {};
\node[text=orange!90!black,font=\fontsize{20}{20}\selectfont] at (1.35,5.83) {\faCogs};
\node[head] at (2.45,6.08) {Workflow scan};
\node[sub] at (2.45,5.62) {GitHub Actions YAML};

\node[card={orangeLine},minimum width=7.55cm,minimum height=1.65cm,anchor=north west] (cmd) at (0.4,4.4) {};
\node[iconbox,fill=orange!15] at (1.35,3.58) {};
\node[text=orange!90!black,font=\fontsize{20}{20}\selectfont] at (1.35,3.58) {\faTerminal};
\node[head] at (2.45,3.83) {Command matching};
\node[sub] at (2.45,3.37) {Regexp per tool};

\node[card={orangeLine},minimum width=7.55cm,minimum height=1.65cm,anchor=north west] (cishare) at (0.4,2.15) {};
\node[iconbox,fill=orange!15] at (1.35,1.33) {};
\node[text=orange!90!black,font=\fontsize{20}{20}\selectfont] at (1.35,1.33) {\faChartBar};
\node[head] at (2.45,1.58) {CI share};
\node[sub] at (2.45,1.12) {Bootstrap 95\% CI};

\draw[flow,orange!90!black] (3.95,4.98) -- (3.95,4.52);
\draw[flow,orange!90!black] (3.95,2.73) -- (3.95,2.27);

% 04 reconstruction
\node[panel={blueLine}{blueBg},minimum width=8.35cm,minimum height=8.25cm,anchor=north west] (history) at (9.40,8.65) {};
\fill[blue] (10.40,7.65) circle (0.62);
\node[white,font=\bfseries\fontsize{17}{17}\selectfont] at (10.40,7.65) {04};
\node[section,text=blue] at (11.40,8.05) {\textls[150]{RECONSTRUCT}};
\node[bigtitle] at (11.40,7.58) {Rule history \& rule rot};

\node[card={blueLine},minimum width=7.55cm,minimum height=1.65cm,anchor=north west] (replay) at (9.80,6.65) {};
\node[iconbox,fill=blue!14] at (10.75,5.83) {};
\node[text=blue!75!black,font=\fontsize{20}{20}\selectfont] at (10.75,5.83) {\faHistory};
\node[head] at (11.85,6.08) {Git history replay};
\node[sub] at (11.85,5.62) {git log -p per file};

\node[card={blueLine},minimum width=7.55cm,minimum height=1.65cm,anchor=north west] (classification) at (9.80,4.4) {};
\node[iconbox,fill=blue!14] at (10.75,3.58) {};
\node[text=blue!75!black,font=\fontsize{20}{20}\selectfont] at (10.75,3.58) {\faThLarge};
\node[head] at (11.85,3.83) {Edit classification};
\node[sub] at (11.85,3.37) {One class per commit};

\node[card={blueLine},minimum width=7.55cm,minimum height=1.65cm,anchor=north west] (survival) at (9.80,2.15) {};
\node[iconbox,fill=blue!14] at (10.75,1.33) {};
\node[text=blue!75!black,font=\fontsize{20}{20}\selectfont] at (10.75,1.33) {\faChartLine};
\node[head] at (11.85,1.58) {Rule rot survival};
\node[sub] at (11.85,1.12) {Kaplan--Meier, 24 months};

\draw[flow,blue!85!black] (replay.south) -- (classification.north);
\draw[flow,blue!85!black] (classification.south) -- (survival.north);

% 05 evaluation
\node[panel={pinkLine}{pinkBg},minimum width=8.35cm,minimum height=8.25cm,anchor=north west] (evaluation) at (18.80,8.65) {};
\fill[pink] (19.80,7.65) circle (0.62);
\node[white,font=\bfseries\fontsize{17}{17}\selectfont] at (19.80,7.65) {05};
\node[section,text=pink] at (20.80,8.05) {\textls[150]{EVALUATE}};
\node[bigtitle] at (20.80,7.58) {ArchUnit adoption};

\node[card={pinkLine},minimum width=7.55cm,minimum height=1.65cm,anchor=north west] (adoption) at (19.20,6.65) {};
\node[iconbox,fill=pink!13] at (20.15,5.83) {};
\node[text=pink,font=\fontsize{20}{20}\selectfont] at (20.15,5.83) {\faCodeBranch};
\node[head] at (21.25,6.08) {Adoption point};
\node[sub] at (21.25,5.62) {First rule commit};

\node[card={pinkLine},minimum width=7.55cm,minimum height=1.65cm,anchor=north west] (eligibility) at (19.20,4.4) {};
\node[iconbox,fill=pink!13] at (20.15,3.58) {};
\node[text=pink,font=\fontsize{20}{20}\selectfont] at (20.15,3.58) {\faCalendarCheck};
\node[head] at (21.25,3.83) {Eligibility};
\node[sub] at (21.25,3.37) {$\leq$12 months history};

\node[card={pinkLine},minimum width=7.55cm,minimum height=1.65cm,anchor=north west] (trend) at (19.20,2.15) {};
\node[iconbox,fill=pink!13] at (20.15,1.33) {};
\node[text=pink,font=\fontsize{20}{20}\selectfont] at (20.15,1.33) {\faChartLine};
\node[head] at (21.25,1.58) {Pre/post trend};
\node[sub] at (21.25,1.12) {Impact replay};

\draw[flow,pink] (adoption.south) -- (eligibility.north);
\draw[flow,pink] (eligibility.south) -- (trend.north);

% One stem leaves discovery, then splits toward workflow and reconstruction.
\coordinate (fork) at (8.45,10.05);
\draw[line width=1.45pt,teal!80!black] (discover.south) -- (fork);
\draw[thickflow,teal!85!black] (fork) -| (workflow.north);
\draw[thickflow,blue!75!black] (fork) -| (history.north);
\draw[thickflow,blue!75!black] ($(discover.north east)+(0,-1.05)$) -- ($(structure.north west)+(0,-1.05)$);
% Links that were missing between the lower panels
\draw[thickflow,teal!85!black] ($(workflow.north east)+(0,-1.85)$) -- ($(history.north west)+(0,-1.85)$);
\draw[thickflow,blue!75!black] ($(history.north east)+(0,-1.85)$) -- ($(evaluation.north west)+(0,-1.85)$);

\end{tikzpicture}%
}
\caption{Study pipeline, from repository discovery and screening through CI detection, rule extraction, history reconstruction, and adoption evaluation.}
\label{fig:study-pipeline}
\end{figure*}

Our methodology comprises five phases, 
as illustrated in Figure~\ref{fig:study-pipeline}.
First, we identify candidate GitHub repositories
and apply eligibility criteria to establish the study corpus. Second,
we analyze GitHub Actions workflows to determine whether architecture
checks are integrated into CI. Third, we extract executable rules from
ArchUnit, import-linter, and dependency-cruiser, normalize their
representations, and classify them by the architectural constraints
they enforce. Fourth, we reconstruct the Git histories of rule files
to identify rule-rot events, defined as the first observed exemption
or disabling of an architectural check. Finally, we assess ArchUnit
adoption by tracking source-level dependency violations over the
12-month periods before and after parseable architecture rules are
introduced.

\subsection{Repository discovery and screening}
We use the GitHub Search API to identify candidate repositories containing
artifacts associated with ArchUnit, import-linter, or dependency-cruiser.
Our search strategy is designed to retrieve repositories that contain
a tool-specific configuration artifact, using a seperate query for
each tool.


After deduplicating the results within each tool,
this process yields 15,542 repository-tool hits, comprising 9,258 for
ArchUnit, 3,108 for import-linter, and 3,176 for dependency-cruiser.

We then exclude forks, archived repositories, template or scaffolding projects, and repositories with less than six months between creation and the latest push.
These exclusions remove 10,799 hits.
Nearly all of them are shorter than six months.
The screen leaves 4,743 eligible repository--tool observations.
Repositories with fewer than two contributors, and repositories without a tag, stay in this set when they pass the screen above.
Of the eligible observations, 422 are unavailable, contain no rule file, or fail to clone.
The comparison corpus is the remainder: 4,321 repository--tool rows from 4,300 repositories, of which 2,422 are ArchUnit, 826 are import-linter, and 1,073 are dependency-cruiser.

\subsection{CI integration detection}
We inspect the default branch of each repository in the comparison corpus and read all \texttt{.yml} and \texttt{.yaml} files under \texttt{.github/}.
This scan covers the 4,321 repository--tool rows. Integration is classified using tool-specific evidence in GitHub Actions workflow text. For import-linter, a workflow is considered integrated if it contains \texttt{lint-imports} or \texttt{import-linter}. For dependency-cruiser, the corresponding indicators are \texttt{depcruise} and \texttt{dependency-cruiser}. ArchUnit does not provide a dedicated command that can be identified in the same way. We therefore use a proxy: a repository is classified as integrated if a workflow invokes Maven \texttt{test}, \texttt{verify}, \texttt{install}, or \texttt{package}, or Gradle \texttt{test}, \texttt{check}, or \texttt{build}. These indicators provide evidence of potentially relevant build or test execution, but do not establish that ArchUnit rules are actually executed.

Integration proportions are calculated over repository--tool observations for which a workflow result can be determined. Unavailable or indeterminate workflow results are excluded from the denominator. We report 95\% confidence intervals obtained by bootstrapping repositories, preserving the repository as the resampling unit.

\subsection{Rule extraction and classification}

The comparison corpus contains 11,098 rule files.
We parse them into a common construct schema.
Each successfully mapped construct is assigned to one of seven categories: dependency, layering, cycle, naming, annotation, visibility, or other.
Table~\ref{tab:rule_taxonomy} defines each category, the constructs that fall under it, and the number of held-out constructs in that category.
Constructs that cannot be mapped remain uncategorized and are excluded from category-proportion calculations.
The first 1,000 eligible repository--tool rows are the development set on which the category map was fixed, and their constructs are excluded from those proportions.
The classification set contains the remaining 21,713 constructs.
We define the structural family as the combination of dependency and layering constraints, calculated over parsed constructs in that set.

\begin{table}[t]
\centering
\caption{Taxonomy used to classify parsed architecture-rule constructs. Proportions use the seven mapped categories only.}
\label{tab:rule_taxonomy}
\footnotesize
\setlength{\tabcolsep}{4pt}
\renewcommand{\arraystretch}{1.12}
\begin{tabular}{@{}>{\columncolor{white}}p{1.55cm}>{\columncolor{gray!15}}p{3.15cm}>{\columncolor{white}}p{2.55cm}@{}}
\toprule
\textbf{Category} & \textbf{Constraint} & \textbf{Constructs} \\
\midrule
Dependency &
Allowed, forbidden, required, or conditional edges between packages or modules &
\texttt{forbidden}, \texttt{allowed}, \texttt{required} \\
Layering &
Ordered layers and the dependencies those layers permit &
\texttt{layers}, \texttt{layering} \\
Cycle &
Freedom from cycles among slices, siblings, or imports &
\texttt{cycle\_freedom}, \texttt{circular} \\
Naming &
Names or declaration sites of packages and types &
\texttt{naming}, \texttt{package\_location} \\
Annotation &
Presence or form of annotations and stereotypes &
\texttt{annotation} \\
Visibility &
Access modifiers and the exposed API surface &
\texttt{visibility} \\
Other &
A check that states none of the constraints above &
\texttt{inheritance}, \texttt{signature} \\
\midrule
Uncate\-gorized &
Not an architectural constraint, or text the parser does not map. Omitted from category proportions &
\texttt{FreezingArchRule},\newline \texttt{@ArchIgnore}, \texttt{unparsed} \\
\bottomrule
\end{tabular}
\end{table}


The mapping and category assignment are performed automatically according to predefined procedures. We did not construct a manually labelled gold standard or validate the assignments against human annotations. Consequently, we do not report inter-coder agreement, and the resulting category distributions should be interpreted as outputs of the automated classification procedure rather than as manually validated labels.

\subsection{Rule-history reconstruction}

We clone repositories using \texttt{--filter=blob:none} without checking out the working tree, then reconstruct the histories of rule files by examining Git diffs and parsing file contents before and after each revision. This process reconstructs 86,339 edits of the 11,098 rule files. Each revision receives one edit class according to the following precedence order: disable, exempt, loosen, tighten, add, and restructure.

A \emph{disable} event includes deleting or commenting out a rule, adding \texttt{@ArchIgnore}, or turning off a rule's severity. An \emph{exempt} event includes introducing a \texttt{FreezingArchRule}, adding an \texttt{ignore\_imports} list, introducing a dependency-cruiser exception, or downgrading severity. We classify a change as \emph{loosen} or \emph{tighten} according to whether the surviving constraint becomes less or more restrictive. An \emph{add} event introduces a new rule. All remaining changes are classified as \emph{restructure}, including file changes that do not alter the parsed rules.

As defined in Section~\ref{sec:background}, we operationalize \emph{rule rot} at the rule-file level as the first observed exemption or disabling event in a file's history. This measure captures an observable change to the enforcement of architectural constraints; it does not establish why the change occurred or whether it was justified. We estimate the cumulative incidence of the first observed rule-rot event over 24 months using Kaplan--Meier analysis and repository-level bootstrap confidence intervals. We also fit a Cox proportional-hazards model using the logarithms of file-level commit count and author count as covariates. Edit classes are assigned automatically and are not validated against human labels.

\subsection{ArchUnit adoption evaluation}

Adoption is the first commit that introduces a parseable ArchUnit rule.
The comparison corpus contains 2,422 ArchUnit repository--tool rows.
A repository enters the trajectory cohort when at least 12 months of its history lie before that commit and 12 months lie after it.
The window retains 317 repositories.

We reconstruct source-level import graphs at historical snapshots and evaluate only dependency and layering constraints that can be assessed from those graphs. This is a replay-based analysis of detectable violations, not an execution of ArchUnit itself. On average, the evaluated constraint family covers 52.9\% of the identified rules, while 33.5\% of all evaluable rules can be assessed by the replay procedure. Repositories without any evaluable rules are excluded from the violation analysis.

For repositories observed at one month before adoption, at adoption, and 12 months afterward, we report violation counts and classify their trajectories as decreasing, unchanged, or increasing. This complete three-point observation set contains 141 repositories. We define a \emph{ratifier} as a repository with zero evaluable violations one month before adoption. A \emph{reformer} is a repository with at least one evaluable violation during the pre-adoption period and zero violations at the adoption point. We estimate within-repository trends using ordinary least squares.

This analysis is descriptive and restricted to repositories that adopt ArchUnit and satisfy the relevant observation criteria. It has no untreated control group and does not employ a difference-in-differences design. Accordingly, changes in violation counts are interpreted as trajectories observed around rule introduction, not as causal effects attributable to ArchUnit adoption.
